
\section{本章内容}

本章是全书的开篇，回答三个问题：本体论从哪里来（哲学源头）、
工程本体论研究什么（形式化、可计算的知识结构）、
为什么 AI 时代它变得不可或缺（大模型的幻觉问题与知识权威源）。

\section{文件说明}

\begin{center}
\begin{tabularx}{\linewidth}{@{}XX@{}}
\toprule
\textbf{文件} & \textbf{内容} \\
\midrule
\texttt{philosophy-to-engineering.txt} & 哲学本体论与工程本体论的对照与转向 \\
\texttt{ontology-in-ai-era.txt} & AI时代的本体定位：幻觉控制与知识权威源 \\
\texttt{book-roadmap.txt} & 全书知识地图与三条学习路径 \\
\bottomrule
\end{tabularx}
\end{center}

\section{核心观点}

\begin{center}
\begin{tabularx}{\linewidth}{@{}lXX@{}}
\toprule
\textbf{维度} & \textbf{哲学本体论} & \textbf{工程本体论} \\
\midrule
研究对象 & 存在的本质、世界的本源 & 领域概念、关系与规则的形式化表示 \\
方法 & 思辨、追问 & 建模、形式化、计算 \\
判据 & 无法收敛（高熵） & 一致性可检验、可推理（低熵） \\
产出 & 哲学体系 & 可被机器处理的知识结构 \\
\bottomrule
\end{tabularx}
\end{center}

\section{经典定义}

\begin{noteblock}
\textbf{Gruber (1993)}: "An ontology is an explicit specification of a conceptualization."
（本体是概念化的显式规范）

\textbf{Studer et al. (1998)} 综合扩展为：本体是\textbf{共享}概念化的\textbf{形式化}、\textbf{显式}规范
—— 四个关键词：概念化、显式、形式化、共享。
\end{noteblock}

\section{学习建议}

\begin{itemize}
\item 零基础读者：按第1\(\rightarrow\)2\(\rightarrow\)3\(\rightarrow\)4章顺序阅读，先建立概念再学语言
\item 有知识图谱经验的读者：可从第3章方法论切入，第2章作为理论参考
\item 关注 AI 应用的读者：第1章后可直接跳读第8章，再回补第4、5章技术细节
\end{itemize}
